\documentclass[11pt]{article}
% Shared preamble for the MSH_Tensor_Network_Methods function references (compile with tectonic / XeLaTeX)
\usepackage{fontspec}
\usepackage{unicode-math}
\setmainfont{texgyrepagella}[Extension=.otf,UprightFont=*-regular,BoldFont=*-bold,ItalicFont=*-italic,BoldItalicFont=*-bolditalic]
\setmathfont{texgyrepagella-math.otf}
\setmonofont{DejaVu Sans Mono}[Scale=0.82]
\newfontfamily\symfb{DejaVu Sans}[Scale=0.82]
\usepackage{newunicodechar}
\newunicodechar{ℓ}{{\symfb ℓ}}
\newunicodechar{∥}{{\symfb ∥}}
\newunicodechar{≪}{{\symfb ≪}}
\usepackage[a4paper,margin=2.2cm]{geometry}
\usepackage{amsmath}
\usepackage{booktabs,tabularx,longtable,array}
\usepackage{graphicx}
\usepackage{xcolor}
\usepackage{float}
\usepackage{caption}
\captionsetup{font=small,labelfont=bf}
\usepackage{fvextra}
\fvset{breaklines,breakanywhere,fontsize=\small,frame=leftline,framerule=1.5pt,rulecolor=\color{gray!45},xleftmargin=6pt}
\usepackage{enumitem}
\setlist{nosep,leftmargin=*}
\usepackage[hidelinks]{hyperref}
\usepackage{titlesec}
\titleformat{\section}{\Large\bfseries}{\thesection}{0.6em}{}
\titleformat{\subsection}{\large\bfseries}{\thesubsection}{0.6em}{}
\titleformat{\subsubsection}{\normalsize\bfseries}{\thesubsubsection}{0.6em}{}
\setlength{\parskip}{0.45em}\setlength{\parindent}{0pt}
% one box per function: \fn{name}{signature} then a description
\newcommand{\fnname}[1]{\texttt{\textbf{#1}}}
\newenvironment{function}[1]{\par\medskip\noindent\textbf{\large\texttt{#1}}\par\nopagebreak}{\par\medskip}
\newcommand{\code}[1]{\texttt{#1}}

\usepackage{mathtools}
\sloppy\emergencystretch=3em
\usepackage{needspace}
\renewenvironment{function}[1]{\par\Needspace{7\baselineskip}\medskip\noindent{\raggedright\textbf{\large\texttt{\hyphenchar\font=-1 #1}}\par}\nopagebreak}{\par\medskip}
% allow line breaks after underscores in \code{}
\renewcommand{\code}[1]{\texttt{\hyphenchar\font=-1 #1}}
\hypersetup{pdftitle={MSH chain on a 2D substrate: function reference}}
\newcommand{\ket}[1]{|#1\rangle}
\newcommand{\bra}[1]{\langle#1|}
\newcommand{\TODO}[1]{\textbf{[TODO: #1]}}

\begin{document}

\begin{center}
{\LARGE\bfseries MSH chain on a 2D substrate}\\[0.5em]
{\Large Function reference}\\[1.2em]
\resizebox{\linewidth}{!}{\begin{tabular}{ll}
\code{msh\_substrate\_tn.ipynb} & Julia, TensorBinding.jl / ITensors: quantics MPO, penalized DMRG\\
\code{msh\_substrate\_python.ipynb} & Python, SciPy / MUMPS: sparse matrix, dense and shift-invert reference\\
\code{msh\_substrate\_scaling.ipynb} & Python: cost of the methods against system size (plots)\\
\end{tabular}}\\[1.2em]
Repository \code{MSH\_Tensor\_Network\_Methods}, branch \code{msh-substrate}\\[0.3em]
30 October 2026
\end{center}

\vspace{1em}
\noindent This document explains every function and struct defined in the three notebooks (folder \code{msh\_substrate/}): what it
computes, the method, its inputs and outputs, the numerical settings that matter, and where it is used. Numbers stated in the
markdown of the notebooks are quoted as such; numbers from the outputs are those of the executed notebooks in the repository. Function names in
\code{typewriter} are those of the code.

\tableofcontents
\clearpage

%======================================================================
\section{Overview}
%======================================================================

The three notebooks measure how the cost of computing the Majorana zero modes of an MSH chain grows when the chain sits on a 2D
superconducting substrate that is made larger and larger, and compare a tensor-network method with dense and sparse-direct methods.
The model is a $2^{L_x}\times2^{L_y}$ square lattice with hopping, Rashba coupling and s-wave pairing on every site, and a centred magnetic
chain of $\ell$ sites in the middle row (section~\ref{sec:model}). Two quantities are computed at every size: the splitting $E_f$ of the
zero-mode pair and its charge difference $\Delta Q = \bra{\psi_f}\tau_z\ket{\psi_f}$.

\begin{itemize}
\item \code{msh\_substrate\_tn.ipynb} (Julia) builds the BdG Hamiltonian as a quantics MPO on $2+L_y+L_x$ qubits, finds the near-zero subspace
  by penalized DMRG on $K = H^2+\eta^2$, localizes the two Majoranas with the position operator and pairs them into $\psi_f$.
\item \code{msh\_substrate\_python.ipynb} (Python) builds the same Hamiltonian as a sparse matrix and solves it by dense \code{eigh} and by
  shift-invert Lanczos with SuperLU or MUMPS; it also checks the Hamiltonian against the MPO, the convergence in the substrate size, the
  charge-difference formula, and scans the parameters for a point where the zero modes leak into the substrate.
\item \code{msh\_substrate\_scaling.ipynb} (Python) reads the CSV files of the other two and compares time, memory, accuracy and bond
  dimensions, and finds the crossover between the methods.
\end{itemize}
Two parameter points are used: the benchmark point $\mu = -7$, $\Delta = 2.4$, $J = 5.8$, $\ell = 12$, where the substrate is a gapped
trivial superconductor and the zero modes stay on the chain row, and the leaking point $\Delta = 0.6$, $\mu = -3$, $J = 3$, $\ell = 64$
(gap 0.090), where they spread over several rows.

\paragraph{Main results.}
\begin{itemize}
\item \emph{Benchmark point.} The bond dimensions do not grow with the system size: $\chi(H) = 7$, $\chi(K)\le38$, $\chi(\psi)\le26$ from one row
  to $2048\times2048$ sites. The TN time grows as $N^{0.14}$ (126 s at $1024^2$, 140 s at $2048^2$, 1.8 GB process memory), against
  $N^{1.5}$ for SuperLU and MUMPS (MUMPS: 70 s at $512^2$, 1013 s and 28.6 GB at $1024^2$; SuperLU runs out of memory there) and $N^{3.07}$ for
  dense diagonalization. The measured crossover is at about $589^2$ sites. The MUMPS times at this point are inflated by subnormal
  floating-point arithmetic (with flush-to-zero the $1024^2$ factorization takes 107 s instead of 771 s); corrected, the crossover is
  at roughly $700^2$. TN errors: $\le2.5\times10^{-8}$ in $\Delta Q$ and $\le1.1\times10^{-11}$ in $E_f$.
\item \emph{Leaking point.} The bond dimensions are larger but still flat in the size ($\chi_\psi = 114$--137, $\chi_K = 31$--37), and the TN
  time is 1.3--3.0 thousand seconds from $256^2$ to $2048^2$. MUMPS is faster at every size up to $1024^2$ (295 s) and about equal at
  $2048^2$ (estimate 2326 s against 2556 s), where it needs $\sim$121 GB against 2.3 GB for the TN. TN errors $\le2.6\times10^{-12}$ in $E_f$ and
  $\le9\times10^{-10}$ in $\Delta Q$.
\item \emph{DMRG details that matter.} $K$ must be built with \code{cutoff\_K = 1e-20} for $\Delta Q$ at the $10^{-9}$ level; the second DMRG state
  is started from the mirror image of the first, which fixes a convergence failure at $2048^2$; and which basis of the zero-mode pair DMRG
  returns (Majoranas, $\chi\approx114$, or two-ended mixtures, $\chi\approx136$) varies with the random start, so the time of a size is uncertain
  by up to a factor 1.7.
\end{itemize}

\subsection{Model and conventions}\label{sec:model}

\paragraph{Lattice.} A square substrate of $N_x\times N_y = 2^{L_x}\times2^{L_y}$ sites with open boundaries in both directions;
$N = N_xN_y$ sites and BdG dimension $D = 4N$. Sites are labelled by the 0-indexed column $i_x$ and row $i_y$. The magnetic chain
is the centred segment of $\ell$ sites in the middle row,
\[
i_{y0} = \lfloor N_y/2\rfloor,\qquad i_x = x_0,\dots,x_0+\ell-1,\quad x_0 = \lfloor (N_x-\ell)/2\rfloor
\]
(\code{chain\_row}, \code{chain\_x} in Julia; \code{chain\_sites} in Python). With $L_y = 0$ there is a single row and the system is the
MSH chain of \code{msh\_chain\_tn.ipynb} with a centred topological segment.

\paragraph{Hamiltonian.} The particle block is
\begin{gather*}
H_0 = \sum_{\mathbf r} c^\dagger_{\mathbf r}(-\mu)c_{\mathbf r}
 + \sum_{\mathbf r}\bigl[c^\dagger_{\mathbf r+\hat x}h_x c_{\mathbf r} + c^\dagger_{\mathbf r+\hat y}h_y c_{\mathbf r} + \text{h.c.}\bigr]
 + J\sum_{\mathbf r\in\text{chain}} c^\dagger_{\mathbf r}\,\mathbf S(\mathbf r)\cdot\boldsymbol\sigma\,c_{\mathbf r},\\
h_x = -t + i\alpha\sigma_y,\qquad h_y = -t - i\alpha\sigma_x ,
\end{gather*}
with $s$-wave singlet pairing $\Delta\,i\sigma_y$ on every site; in the Julia Nambu basis
\[
H_{\rm BdG} = \begin{pmatrix} H_0 & \Delta\,i\sigma_y\\ (\Delta\,i\sigma_y)^\dagger & -H_0^*\end{pmatrix}.
\]
Two spin textures are used: \code{S\_z}, all moments along $\hat z$ (topological; every benchmark), and \code{S\_tilt},
$\mathbf S(i_x) = (\sin 0.3i_x\cos0.7,\ \sin0.3i_x\sin0.7,\ \cos0.3i_x)$, a generic texture with all three components, used only to check
the Hamiltonian. The texture depends on $i_x$ only, since the chain lies in one row.

\paragraph{Parameter points.}
\begin{center}
\resizebox{\linewidth}{!}{\begin{tabular}{lllll}
\toprule
Symbol & Benchmark point & Second point & Code name (Julia / Python) & Meaning\\
\midrule
$t$ & 1 & 1 & \code{t} / \code{t} & nearest-neighbour hopping\\
$\mu$ & $-7$ & $-3$ & \code{μ} / \code{mu} & chemical potential\\
$\alpha$ & 0.9 & 0.9 & \code{α} / \code{alpha} & Rashba spin--orbit coupling\\
$\Delta$ & 2.4 & 0.6 & \code{Δ} / \code{Delta} & $s$-wave pairing\\
$J$ & 5.8 & 3 & \code{J} / \code{J} & exchange coupling to the moments\\
$\ell$ & 12 & 64 & \code{ℓ} / \code{ell} & chain length (zero-mode runs)\\
\midrule
used in & TN 7, Py 5--7 & TN 8, Py 8 & & \\
\bottomrule
\end{tabular}}
\end{center}
The size scans without eigensolver use $\ell = 48$ (width scan, and Python section 3) and $\ell = 3N_x/4$ (square scan).

The markdown describes the two points as follows. At the benchmark point the bare band ($\pm4.74$ with the Rashba term) is empty,
so the substrate is a gapped, trivial superconductor on which the chain stays topological; the zero modes barely enter it (gap
$\approx0.55$, decay within one row; only $\sim2\times10^{-5}$ of the weight lies four or more rows away), which favours the tensor network.
With $\ell = 12$ the pair splitting $E_f\approx4\times10^{-4}$ and the charge difference $|\Delta Q|\approx3\times10^{-3}$ are far above
round-off, so they measure accuracy ($\ell = 48$ would give a splitting $\sim10^{-14}$). The second point, found by the
$\mu\times J\times\Delta$ scan of Python section 8, is inside the band, topological and gapped (gap $0.0902$), and its zero mode spreads over
several rows: 0.32 of its weight on the chain row, 0.13 four or more rows away, decaying as $e^{-|\Delta y|/2}$. The splitting decays
along the chain with a length of about 7 sites; $\ell = 64$ gives $E_f/E_{\rm gap} = 10^{-4}$, comparable with the benchmark point. The
converged exact values quoted there are $E_f = 9.04771\times10^{-6}$ and $\Delta Q = 4.58344\times10^{-4}$.

\paragraph{Quantities.} Every zero-mode run returns three numbers: the pair energy $E_f>0$ of the fermion
$\psi_f = (\gamma_1+i\gamma_2)/\sqrt2$ built from the two Majoranas, the charge difference
\[
\Delta Q = \bra{\psi_f}\tau_z\ket{\psi_f} = \sum_{\mathbf r,\sigma}\bigl(|u_\sigma(\mathbf r)|^2 - |v_\sigma(\mathbf r)|^2\bigr),
\]
which equals the charge difference of the two logical states (Python section 7), and the gap $E_{\rm gap}$, the next level above the pair.

\subsection{Quantics encoding (Julia)}

A single-particle BdG vector is an MPS, and $H$ an MPO, on $L_x+L_y+2 = \log_2N+2$ sites of dimension 2:
\[
[\,\text{nambu};\ \text{spin};\ y_1,\dots,y_{L_y};\ x_1,\dots,x_{L_x}\,],
\]
most significant bit first. The column and row are $i_x = \sum_{k=1}^{L_x} b^x_k2^{L_x-k}$ and $i_y = \sum_{k=1}^{L_y} b^y_k2^{L_y-k}$, so the
position register (y bits, then x bits) read as one binary number is
\[
n = i_x + i_yN_x,
\]
TensorBinding's 2D row-major convention. A shift by one site in $x$ is therefore $\ket n\to\ket{n+1}$ and a shift by one row is
$\ket n\to\ket{n+N_x}$, both binary adders on the whole register. In the dense form of the MPO (\code{dense}, big-endian over the
sites) the index is $n_b\cdot2N + \sigma N + n$ with $n_b,\sigma\in\{0,1\}$. The x qubits are the sites
$3+L_y,\dots,2+L_y+L_x$ (\code{xqubits}).

\subsection{Nambu bases of the two codes}

\begin{center}
\begin{tabular}{lll}
\toprule
 & Julia (TensorBinding) & Python (kwant convention)\\
\midrule
Spinor per site & $(c_\uparrow, c_\downarrow, c^\dagger_\uparrow, c^\dagger_\downarrow)$ & $(c_\uparrow, c_\downarrow, c^\dagger_\downarrow, -c^\dagger_\uparrow)$\\
Matrix index & $n_b\cdot2N + \sigma N + n$ & $4n + c$\\
On-site term & $-\mu + J\mathbf S\cdot\boldsymbol\sigma$; pairing $\Delta\,i\sigma_y$ & $-\mu\tau_z + \Delta\tau_x + J\,\tau_0\mathbf S\cdot\boldsymbol\sigma$\\
Hopping $H_{\mathbf r+\hat d,\mathbf r}$ ($\hat x$, $\hat y$) & $-t + i\alpha\sigma_y$,\ $-t - i\alpha\sigma_x$ & $-t\tau_z + i\alpha\tau_z\sigma_y$,\ $-t\tau_z - i\alpha\tau_z\sigma_x$\\
Particle--hole operator & $C = \tau_xK$ (\code{Cconj}) & $C = \tau_y\sigma_yK$\\
Charge operator & $\tau_z$ on the Nambu site (\code{τz\_on}) & $+1$ on $c = 0,1$, $-1$ on $c = 2,3$\\
\bottomrule
\end{tabular}
\end{center}
Per site, the Python components $(a,b,c,d)$ are $(u_\uparrow, u_\downarrow, v_\downarrow, -v_\uparrow)$ in terms of the Julia components
$(u_\uparrow,u_\downarrow,v_\uparrow,v_\downarrow)$. The signed permutation \code{julia\_to\_python} (Python section 2) implements this map; $\tau_z$
and therefore $\Delta Q$ are the same in both bases.

\subsection{Workflow and data files}

Run order (all paths relative to \code{msh\_substrate/}):
\begin{center}\code{msh\_substrate\_tn.ipynb} $\to$ \code{msh\_substrate\_python.ipynb} $\to$ \code{msh\_substrate\_scaling.ipynb}.\end{center}
Python section 2 needs \code{data/tn\_check\_H.npz} from the TN notebook; Python section 8c needs \code{py\_mumps\_estimates.csv} from
Python section 6; the scaling notebook only loads files. All timings are single-threaded (\code{BLAS.set\_num\_threads(1)} in Julia;
\code{OMP\_NUM\_THREADS}, \code{MKL\_NUM\_THREADS}, \code{OPENBLAS\_NUM\_THREADS} set to 1 before NumPy is imported), except the parameter scan of
Python section 8a, which uses 8 processes. The markdown gives about 2.5 hours for TN section 8; from the recorded times, TN section 7 takes about 12 minutes and section 8 about 2.5 hours (8924 s); Python section 5 takes at least 40 minutes (dense $64\times64$: 886 s, MUMPS $1024^2$: 1013 s) and section 8 several minutes (scan 198 s, MUMPS $1024^2$: 295 s); the scaling notebook runs in seconds.

\begin{center}
\small
\begin{tabularx}{\textwidth}{>{\raggedright}p{4.6cm}>{\raggedright}p{3.6cm}>{\raggedright\arraybackslash}X}
\toprule
File & Written by & Content\\
\midrule
\code{data/tn\_check\_H.npz} & TN sec.~4 & ten dense MPO matrices, keys \code{H\_Lx\_Ly\_ℓ\_\{z,tilt\}} (not committed)\\
\code{data/zm\_progress.log} & TN sec.~7 & one line per finished size of the benchmark\\
\code{tn\_zero\_modes.csv} & TN sec.~7 & TN benchmark, 22 columns (committed)\\
\code{tn\_candidate.csv} & TN sec.~8 & TN at the second point, 24 columns (committed, resumable)\\
\code{py\_zero\_modes.csv} & Py sec.~5 & dense / SuperLU / MUMPS benchmark (resumable)\\
\code{py\_mumps\_estimates.csv} & Py sec.~6 & MUMPS analysis estimates, $L = 7\dots11$\\
\code{py\_mu\_scan.csv} & Py sec.~8a & 1305 points of the $\mu\times J\times\Delta$ scan\\
\code{py\_candidate.csv} & Py sec.~8c & MUMPS at the second point (resumable)\\
\bottomrule
\end{tabularx}
\end{center}

\paragraph{Mapping between the codes.}
\begin{center}
\small
\begin{tabularx}{\textwidth}{>{\raggedright}p{3.2cm}>{\raggedright}p{5.4cm}>{\raggedright\arraybackslash}X}
\toprule
Quantity & Julia (\code{msh\_substrate\_tn}) & Python (\code{msh\_substrate\_python})\\
\midrule
Geometry & \code{Substrate}, \code{Nx}, \code{Ny}, \code{chain\_row}, \code{chain\_x} & \code{chain\_sites}\\
Texture & \code{S\_z}, \code{S\_tilt} & \code{S\_z}, \code{S\_tilt}\\
$H$ & \code{msh\_substrate} (MPO) & \code{ham} (CSR)\\
Check of $H$ & section 4 (export) & \code{julia\_to\_python}, section 2\\
Size of $H$ & \code{mpo\_size}, section 5 & \code{size\_info}, section 4\\
Zero modes & \code{zero\_mode\_subspace}, \code{fermion\_from\_subspace} & \code{solve}, \code{pair\_values}\\
$\Delta Q$ & \code{ΔQ\_mode} & \code{pair\_values}; \code{mask\_check} (section 7)\\
Benchmark & \code{zm\_run} (section 7) & \code{solve} in \code{in\_fresh\_process} (section 5)\\
Second point & \code{cand\_run} (section 8) & \code{solve\_candidate} (section 8c)\\
\bottomrule
\end{tabularx}
\end{center}

\clearpage
%======================================================================
\section{Part I --- Tensor-network notebook (\code{msh\_substrate\_tn.ipynb})}
%======================================================================

The Hamiltonian is a quantics MPO on $\log_2N+2$ sites; the zero modes come from penalized DMRG on $K = H^2+\eta^2$. Dense matrices
appear only in section 4, to export check data for the Python notebook. Packages: \code{TensorBinding} (alias \code{TB}), \code{ITensors},
\code{ITensorMPS}, \code{LinearAlgebra}, \code{Printf}, \code{NPZ}, \code{Random}, \code{Logging}, \code{DelimitedFiles}.

\subsection{Section 1: setup and helpers}

Global settings: \code{BLAS.set\_num\_threads(1)} (single thread, as in the Python notebook, so that timings compare like with like);
the Pauli matrices \code{σ0, σx, σy, σz} (\code{ComplexF64}); \code{DATA\_DIR = "data"}, created if missing, for the check data.

\begin{function}{diag\_mpo(ψ, pos)}
\begin{Verbatim}
δ = ITensor(ComplexF64, s, s', dag(s''))
for j in 1:2; δ[s => j, s' => j, s'' => j] = 1; end
W[k] = replaceprime(ψ[k] * δ, 2 => 0; tags=tags(s))
\end{Verbatim}
Turns an MPS $\psi$ on the position qubits \code{pos} into the diagonal MPO $\mathrm{diag}(\psi) = \sum_n\psi_n\ket n\bra n$. Each site tensor
is contracted with the copy tensor $\delta_{ss's''}$, which places the value $\psi[k]_{s=j}$ on the diagonal $(s', s) = (j, j)$; the
\code{replaceprime} renames $s''$ back to $s$. The bond dimension of the MPO equals that of $\psi$, and the operation is exact. Used for
the x-bond mask and the exchange in \code{msh\_substrate}, and by \code{diag\_mpo\_exact}.
\end{function}

\begin{function}{vector\_mps(v, q; cutoff=1e-14)}
MPS of a vector $v$ of length $2^{|q|}$ on the qubits \code{q} (\code{q[1]} the most significant bit). The vector is reshaped into a
$|q|$-index array, the indices reversed with \code{permutedims} (Julia's column-major reshape puts the least significant bit first), and
compressed by TT-SVD (\code{MPS(A, q; cutoff)}) with truncation \code{cutoff = 1e-14}, i.e.\ exact up to round-off.
\code{@disable\_warn\_order} suppresses ITensors' warning about tensors with many indices. Cost and memory are $O(2^{|q|})$, since the full
array is built.
\end{function}

\begin{function}{diag\_mpo\_exact(v, pos; cutoff=1e-14)}
\code{diag\_mpo(vector\_mps(v, pos; cutoff), pos)}: the exact diagonal MPO $\sum_nv_n\ket n\bra n$ of a full-length vector. Not used by the
current Hamiltonian; an earlier version compressed length-$N$ vectors with it, which, according to the markdown of section 3, took 12--14 s
of the build at $512\times512$ and $1024\times1024$ sites. It is kept as the reference construction.
\end{function}

\begin{function}{row\_profile\_mps(yvecs, w, ypos, xpos; cutoff=1e-14)}
\begin{Verbatim}
ψx = vector_mps(w, xpos; cutoff)
isempty(ypos) && return ψx
link = [Index(1, "Link,ly=$k") for k in eachindex(ypos)]
...  t = ITensor(ComplexF64.(yvecs[k]), s) * onehot(link[k] => 1)
ψx[1] *= onehot(link[end] => 1)
return MPS(vcat(T, collect(ψx)))
\end{Verbatim}
The MPS of $(\mathbf y_1\otimes\dots\otimes\mathbf y_{L_y})\otimes w$ on $[\,y_1..y_{L_y};\ x_1..x_{L_x}\,]$: a product of two-component vectors
$\mathbf y_k$ on the y qubits times a profile $w$ of length $N_x$ on the x qubits. Only $w$ is compressed (\code{vector\_mps}); each y qubit gets a
single tensor joined to its neighbours by dimension-1 links, and the last link is attached to the first x tensor. The cost is
$O(N_x)$ instead of $O(N)$. Two choices of \code{yvecs} are used: the bits of one row $i_y$ (each $\mathbf y_k = (1,0)$ or $(0,1)$), which gives a
vector that is $w(i_x)$ on that row and zero elsewhere; and $\mathbf y_k = (1,1)$ for all $k$, which gives the same profile $w$ on every row.
With $L_y = 0$ it returns $\psi_x$ unchanged. This is what makes the build of $H$ cost $O(N_x)$ (section 3).
\end{function}

\begin{function}{dense(mpo, sites), adjoint\_mpo(K)}
\code{dense} contracts the whole MPO (\code{prod}) and reshapes it into a $D\times D$ matrix, big-endian over \code{sites} (first site most
significant), giving the index $n_b\cdot2N+\sigma N+n$. Small systems only (section 4). \code{adjoint\_mpo(K) = swapprime(dag(K), 0, 1)} is
the Hermitian conjugate of an MPO: complex conjugation of the tensors and exchange of bra and ket indices.
\end{function}

\subsection{Section 2: geometry}

\begin{function}{struct Substrate; Nx, Ny, nsites, chain\_row, chain\_x}
\begin{Verbatim}
struct Substrate
    Lx::Int; Ly::Int; ℓ::Int
end
chain_row(g::Substrate) = Ny(g) ÷ 2
chain_x(g::Substrate) = ((Nx(g) - g.ℓ) ÷ 2):((Nx(g) - g.ℓ) ÷ 2 + g.ℓ - 1)
\end{Verbatim}
\code{Substrate(Lx, Ly, ℓ)} holds the geometry: $N_x = 2^{L_x}$ columns (\code{Nx}), $N_y = 2^{L_y}$ rows (\code{Ny}), $N$ sites (\code{nsites}) and
a chain of $\ell$ sites. \code{chain\_row} is the chain row $\lfloor N_y/2\rfloor$ (0 for one row), \code{chain\_x} the 0-indexed range of chain
columns. For even $\ell$ (all zero-mode runs) the chain is exactly centred, $x\to N_x-1-x$ maps it onto itself; for odd $\ell$ (one check case)
it is off by half a site.
\end{function}

\begin{function}{S\_z(ix), S\_tilt(ix)}
Spin textures as functions of the column: \code{S\_z} returns $(0,0,1)$; \code{S\_tilt} returns
$(\sin0.3i_x\cos0.7,\ \sin0.3i_x\sin0.7,\ \cos0.3i_x)$, so that all three exchange components are nonzero and vary along the chain. Identical
lambdas exist in the Python notebook.
\end{function}

\subsection{Section 3: Hamiltonian MPO}

\begin{function}{msh\_substrate(g, S; t=1.0, μ=-7.0, α=0.9, Δ=2.4, J=5.8, cutoff=1e-13, sites=nothing)}
Returns \code{(HB, [ns; ss; pos])}: the BdG MPO of the chain with texture \code{S(ix)} on substrate \code{g}, and its site list. \code{ns} and
\code{ss} are TensorBinding's Nambu and spin indices (\code{TB.nambu\_index()}, \code{TB.spin\_index()}), \code{pos} the $L_x+L_y$ qubits
(\code{siteinds("Qubit", ...)}), split into \code{ypos} (first $L_y$) and \code{xpos}. If \code{sites} is given its indices are reused.
\code{P(op, M) = TB.prepend\_spin(M, ss, op)} puts a $2\times2$ spin matrix in front of a position MPO.

\emph{Hopping.} The x bonds and the y bonds are built from TensorBinding's binary-adder MPO \code{TB.shift\_pair\_mpos(pos, q; cyclic=false)},
whose first element maps $\ket n\to\ket{n+q}$ on the whole register (bond dimension 2: the carry). With \code{cyclic=false} a carry out of the
most significant bit is dropped, so states with $n+q\ge N$ are mapped to zero.
\begin{itemize}
\item \textbf{x bonds:} $K_x = K_1\cdot\mathrm{mask}$, with $K_1 = $ shift by 1 and
  \[
  \mathrm{mask} = \mathrm{diag}\bigl((1,1)^{\otimes L_y}\otimes m\bigr),\qquad m(i_x) = \begin{cases}0 & i_x = N_x-1\\ 1 & \text{otherwise.}\end{cases}
  \]
  Without the mask the shift would join the last column of row $i_y$ to the first column of row $i_y+1$. The mask is built by
  \code{row\_profile\_mps(every\_row, m, ypos, xpos)} and \code{diag\_mpo}, the product by \code{apply(K1, mask; cutoff)} (the mask acts first).
  Spin matrix $h_x = -t\sigma_0 + i\alpha\sigma_y$.
\item \textbf{y bonds:} $K_y = $ shift by $N_x$ (\code{shift\_pair\_mpos(pos, nx)}), which adds 1 to $i_y$ and maps the top row to zero (open
  boundary); $h_y = -t\sigma_0 - i\alpha\sigma_x$. Only added when $L_y>0$: with one row, $N_x = 2^{L}$ has no bit inside the register, the adder
  keeps only the low bits of $q$ and would act as the identity.
\end{itemize}
For each bond the terms $h\otimes K$ and $h^\dagger\otimes K^\dagger$ (\code{P(Matrix(hop'), adjoint\_mpo(K))}) are added, i.e.\
$c^\dagger_{\mathbf r+\hat d}hc_{\mathbf r}$ and its Hermitian conjugate.

\emph{Exchange.} For each spin component $k$ the vector $w(i_x) = JS_k(i_x)$ on the chain columns (zero elsewhere) becomes
\code{diag\_mpo(row\_profile\_mps(row(iy0), w, ypos, xpos), pos)}: nonzero on the chain row only, where \code{row(iy)} lists the bits of $i_y$
as unit vectors, most significant first. The term $\sigma_k\otimes(\cdot)$ is added only when $\max|w|>10^{-14}$, so \code{S\_z} gives a single
exchange term.

\emph{Assembly.} $H_0 = -\mu\,\sigma_0\otimes\mathbb 1 + \sum_{\rm bonds} + \sum_k$ exchange, summed with \code{+(terms...; cutoff)} at
\code{cutoff = 1e-13}. The pairing MPO is \code{Hp = P(:iSy, Δ * Id)} $=\Delta\,i\sigma_y\otimes\mathbb1$. The BdG MPO is
\[
H_B = \tau_{pp}\otimes H_0 + \tau_{hh}\otimes(-H_0^*) + \tau_+\otimes H_p + \tau_-\otimes H_p^\dagger,
\]
with \code{TB.prepend\_nambu} and the operators \code{:Pp}, \code{:Ph}, \code{:tp} $=\ket p\bra h$, \code{:tm} $=\ket h\bra p$; $-H_0^*$ is
\code{-1.0 * dag(H0)} (\code{dag} conjugates the elements) and $H_p^\dagger$ is \code{adjoint\_mpo(Hp)}. Again summed at \code{cutoff = 1e-13}.

\emph{Cost.} Every diagonal is a product state on the y qubits times a compressed vector of length $N_x$, and the shifts are exact
bond-dimension-2 MPOs, so the build costs $O(N_x)$ (markdown: milliseconds at every size). The resulting $\chi(H)$ is
$\chi(H) = 7$ for spins along $\hat z$ at every size, from one row to $2048\times2048$ sites, and 8--9 for the tilted texture of the check cases; it does not depend on $L_y$.

\emph{Used in:} section 4 (both textures), section 5 (build scan), \code{zm\_run} and \code{cand\_run} (with \code{PC...} for the second point).
\end{function}

\subsection{Section 4: check data for the Python notebook}

For the five geometries \code{check\_cases} $=(L_x,L_y,\ell) = (3,0,6), (3,1,6), (3,2,5), (2,3,3), (4,2,10)$ (one row; two to eight rows; a
substrate taller than wide) and both textures (\code{"z"}, \code{"tilt"}), the cell builds \code{msh\_substrate}, converts it with \code{dense}, and
stores the matrix under the key \code{H\_\$(Lx)\_\$(Ly)\_\$(ℓ)\_\$name} in a dictionary written by \code{npzwrite} to
\code{data/tn\_check\_H.npz}. Dimensions $D = 32$ to 256. For each case it prints $D$, $\chi(H)$ and the Hermiticity error
$\lVert H-H^\dagger\rVert$ ($\chi(H) = 7$ for \code{S\_z}, 8--9 for \code{S\_tilt}; $\lVert H-H^\dagger\rVert = 1.4\times10^{-13}$ to $9.0\times10^{-13}$). The comparison itself is done in
Python section 2.

\subsection{Section 5: MPO size and build time against system size}

\begin{function}{mpo\_size(H)}
\code{sum(length(storage(T)) for T in H)}: the number of stored elements of the MPO tensors, the memory measure compared with the number of
nonzeros of the sparse matrix (Python section 4). Also used for $K$ in \code{zm\_run} and \code{cand\_run}.
\end{function}

\paragraph{Build scan.} \code{scan} is the width scan $(6, L_y, 48)$, $L_y = 0\dots6$, followed by the square scan $(L, L, 3\cdot2^L/4)$,
$L = 4\dots11$ (up to $2048\times2048$ sites). For each size the cell builds $H$ once as a warm-up (so compilation is not timed) and then times
a second build with \code{@elapsed}; it collects $(L_x, L_y, \ell, N, \chi, \text{mem}, t_{\rm build})$ in \code{build} (not written to a file)
and prints $N$, $D$, $\chi(H)$, the MPO elements and the build time. At $2048\times2048$: $\chi(H) = 7$, 3116 MPO elements, build 0.04 s; all builds take 0.01--0.04 s, and the number of elements grows linearly in the number of qubits (612 for one row of 64 sites).

\subsection{Section 6: zero-mode pipeline}

The pipeline of \code{msh\_chain\_tn.ipynb} (sections 8--9) adapted to the substrate: penalized DMRG for the near-zero subspace, $\hat X$
localization and pairing, then $\Delta Q$.

\begin{function}{quietly(f)}
Runs the closure \code{f} with a logger that shows errors only. KrylovKit warns (``ignoring imaginary component'') at every DMRG step on
these complex MPOs; the comment states that the imaginary parts are round-off. Wraps the DMRG calls of \code{zero\_mode\_subspace}.
\end{function}

\begin{function}{norm\_Hψ(H, ψ)}
$\lVert H\psi\rVert$, from the exact product \code{apply(H, ψ; alg="naive", truncate=false)}, orthogonalized to site 1 so that the norm is the
norm of the first tensor. For a normalized eigenvector this is $|E|$; for an approximate one it is an upper bound on the distance of the
nearest eigenvalue from zero. Gives \code{E\_sub} and \code{E\_gap}.
\end{function}

\begin{function}{apply\_site(op, ψ, j), τz\_on(ψ), ΔQ\_mode(ψ)}
\code{apply\_site} applies a single-site operator \code{op} (indices $s'$, $s$) to site $j$ of a copy of $\psi$; bonds are unchanged.
\code{τz\_on} does the same with $\tau_z = \mathrm{diag}(1,-1)$ on the Nambu site (site 1). \code{ΔQ\_mode(ψ) = real(inner(ψ, τz\_on(ψ)))}
$=\bra\psi\tau_z\ket\psi$, the charge difference of a normalized fermion mode.
\end{function}

\begin{function}{mps\_size(ψ), xqubits(g)}
\code{mps\_size}: stored elements of the MPS tensors. \code{xqubits(g) = (3 + g.Ly):(2 + g.Ly + g.Lx)}: site numbers of the x qubits, most
significant first (after Nambu, spin and the $L_y$ y qubits).
\end{function}

\begin{function}{Cconj(ψ)}
Particle--hole conjugate $C\psi = \tau_xK\psi$ of an MPS: \code{dag(ψ)} conjugates all tensors, then $\tau_x$ swaps the particle and hole
components on the Nambu site. The near-zero subspace is closed under $C$.
\end{function}

\begin{function}{mirror\_x(ψ, sites, xq)}
The mirror image $x\to N_x-1-x$ of an MPS. In the quantics encoding $N_x-1-i_x$ is the bitwise complement of $i_x$, so the mirror is a NOT
gate $\begin{psmallmatrix}0&1\\1&0\end{psmallmatrix}$ on every x qubit (site numbers \code{xq}); it is a product operator and leaves the bond
dimensions unchanged. For a centred chain (even $\ell$) it maps one chain end onto the other. Rows and the Nambu and spin components are
untouched.
\end{function}

\begin{function}{zero\_mode\_subspace(HB, sites; nzero=2, η=1e-2, seed=1234, nsweeps=20, ..., xq=nothing)}
Penalized DMRG for \code{nzero} near-zero states and one more state as a gap check. Steps:
\begin{enumerate}
\item $K = H^2+\eta^2$, $\eta = 10^{-2}$, by \code{TB.build\_K(HB, sites, 0.0, η; cutoff\_K)}: $H\cdot H$ by \code{apply} with truncation
  \code{cutoff\_K} and TensorBinding's default cap \code{maxdim\_K = 200}, plus $\eta^2\mathbb1$. Timed as \code{t\_K}.
\item \code{Random.seed!(seed)}; for $k = 1,\dots,$\code{nzero}$+1$ a DMRG run on $K$ from \code{random\_mps(sites; linkdims=10)}. The first state is
  the ground state of $K$, i.e.\ the eigenstate of $H$ closest to zero; later states use \code{dmrg(K, ψs, ψ0; weight)}, which adds the
  penalty $w\sum_j\ket{\psi_j}\bra{\psi_j}$ ($w = $\code{weight}$=10$) for the states already found. All three runs are timed together as
  \code{t\_dmrg}.
\item Each state is normalized. The last one is the gap state $\psi_{\rm gap}$; the first \code{nzero} span the zero-mode subspace.
\end{enumerate}
DMRG settings: \code{nsweeps = 20}, \code{maxdim = [10, 20, 50, 100, 200]} (the last value repeats), \code{cutoff = 1e-14},
\code{noise = [1e-6, 1e-7, 1e-8, 0.0]}, \code{eigsolve\_krylovdim = 20}, \code{eigsolve\_maxiter = 3}, \code{outputlevel = 0}.

\emph{Truncation of $K$} (\code{cutoff\_K = 1e-20}). The markdown explains the choice: $\Delta Q$ is first order in any admixture of bulk
states in the pair (energies only second order), so it is the most sensitive quantity. With \code{cutoff\_K = 1e-14} the one-row chain had
$|\Delta Q_{\rm TN}-\Delta Q_{\rm exact}| = 1.3\times10^{-5}$, from a weight $4.7\times10^{-10}$ outside the exact pair that did not depend on
sweeps or the DMRG cutoff; with \code{1e-20} the error was $1.3\times10^{-10}$. With a substrate the error was a few $\times10^{-9}$ for either value,
limited by DMRG convergence. The tighter cutoff doubles $\chi(K)$ (18 to 38) and makes DMRG 1.2--1.6 times slower on the large substrates.
The inverse-iteration refinement of \code{msh\_chain\_tn.ipynb} is not needed ($E_f\gg2\times10^{-6}$) and was not an alternative: with
$E_f$ close to its shift $\sigma = 10^{-3}$ it made $L_y = 1$ worse ($5.7\times10^{-5}$).

\emph{Start of the second state.} With a random start the second DMRG can end in a wrong state: the markdown reports that at the second
parameter point and $2048\times2048$ sites it stopped at $\lVert H\psi_2\rVert = 0.60$ (a substrate band-edge state, bond dimension at the
cap of 200), and the gap check then picked up the missing zero mode. The second state is therefore started from a state built from
$\psi_1$, after removing its $\psi_1$ component (\code{orth(φ) = +(φ, -inner(ψs[1], φ) * ψs[1]; cutoff=1e-14)}), in this order:
\begin{itemize}
\item if \code{xq} is given, the mirror image \code{mirror\_x(ψ1, sites, xq)}, if its orthogonal part has norm $>0.1$. When $\psi_1$ is a
  Majorana at one chain end this is the Majorana at the other end. Orthogonalized, it keeps norm 0.98--1.00;
\item otherwise $C\psi_1$ (\code{Cconj}), if its orthogonal part has norm $>0.1$ (the subspace is closed under $C$). It is useless when
  $\psi_1$ is a single Majorana: then $C\psi_1\propto\psi_1$ and the orthogonal part has norm $0.000$;
\item otherwise the random MPS (\code{start2 = "random"}).
\end{itemize}
The random MPS of state 2 is drawn in any case, so the random stream of state 3 does not depend on the choice. Both norms are computed
and returned whichever start is taken; \code{C\_perp} $\approx0$ means that $\psi_1$ came out a single Majorana.

\emph{Which basis of the pair DMRG returns.} The near-zero subspace is two-dimensional, and DMRG can return any orthonormal basis of it:
two localized Majoranas, or two-ended mixtures like $\psi_f$. For the exact modes at the second parameter point (TT-SVD with cutoff
$10^{-14}$) a Majorana needs $\chi = 121$ and $\psi_f$ needs $139$, independent of the system size and of the position of the chain on the binary
grid. Which basis comes out is set by the random start of $\psi_1$ and varies from size to size (section 8: $\chi\approx114$ where $\psi_1$ is
a Majorana, $\approx136$ where it is a mixture), so the DMRG time of a size is uncertain by up to $(137/114)^3\approx1.7$. Starting $\psi_1$
from a random MPS projected onto the left half of the substrate was tried and did not fix this at $512^2$ and above (markdown).

\emph{Returns} the named tuple \code{(ψs, E\_sub, E\_gap, K, t\_K, t\_dmrg, start2)}: the \code{nzero} states, their $\lVert H\psi\rVert$, the gap
estimate $E_{\rm gap} = \lVert H\psi_{\rm gap}\rVert$ (which should sit at the bulk gap), $K$, the two times, and
\code{start2 = (kind, C\_perp, mirror\_perp)} with \code{kind} $\in\{$\code{"Cψ1"}, \code{"mirror"}, \code{"random"}$\}$ and the norms of the two
orthogonalized candidates (\code{NaN} where not computed).
\end{function}

\begin{function}{fermion\_from\_subspace(HB, sites, ψs, g; cutoff=1e-14)}
Builds the Majoranas and the fermion mode from the two subspace states.
\begin{enumerate}
\item Overlap matrix $S_{ab} = \langle\psi_a|\psi_b\rangle$ and the position matrix $X_{ab} = \langle\psi_a|\hat X|\psi_b\rangle$ with
  $\hat X = \sum_k2^{L_x-k}\hat n_k$ on the x qubits only ($\hat n_k = \ket1\bra1$ on qubit $k$, applied by \code{apply\_site}); $\hat X$ is the
  column $i_x$, independent of the row.
\item The generalized Hermitian eigenproblem $X c = x\,S c$ (\code{eigen(Hermitian(X), Hermitian(S))}, both symmetrized) gives the Majorana
  centres $x_1\le x_2$ (0-indexed columns) and the coefficients; \code{combine} forms $\gamma_i = c_1\psi_1+c_2\psi_2$ (sum at \code{cutoff}) and
  normalizes it.
\item Phase fix: if $C\gamma = e^{i\theta}\gamma$ then $e^{i\theta/2}\gamma$ is self-conjugate; $\theta = \arg\langle\gamma|C\gamma\rangle$.
\item $\psi_f = (\gamma_1+i\gamma_2)/\sqrt2$, normalized; $E_f = \mathrm{Re}\,\bra{\psi_f}H\ket{\psi_f}$. If $E_f<0$, $\psi_f\to C\psi_f$ and
  $E_f\to-E_f$, so $\psi_f$ is always the $+E$ member.
\end{enumerate}
Returns \code{(centers, γs, ψf, E\_f, C)}. The function computes its own x-qubit range and its own $C$ (the same as \code{xqubits} and
\code{Cconj}).
\end{function}

\subsection{Section 7: zero-mode benchmark}

\code{RUN\_ZM = true}; \code{zm\_file = "tn\_zero\_modes.csv"}; \code{zm\_cols}, the 22 column names. \code{zm\_scan} is the width scan $(6, L_y)$,
$L_y = 0\dots6$, and the square scan $(L,L)$, $L = 4, 5, 7, \dots, 11$ ($L = 6$ is in the width scan), 14 sizes sorted by
$(L_x+L_y, L_y)$, i.e.\ increasing $N$. All with $\ell = 12$ and \code{S\_z}.

\begin{function}{zm\_run(Lx, Ly; ℓ=12)}
One benchmark size: \code{msh\_substrate} (timed as \code{t\_H}), \code{zero\_mode\_subspace} with \code{xq = xqubits(g)},
\code{fermion\_from\_subspace} (timed as \code{t\_pair}) and \code{ΔQ\_mode}. Returns the row
\begin{center}\small
\code{Lx, Ly, ell, N, D, t\_H, t\_K, t\_dmrg, t\_pair, t\_total, chi\_H, chi\_K, chi\_psi, mem\_K\_MB, mem\_psi\_MB, rss\_MB, E\_f, dQ, E\_gap, x1, x2, start2}
\end{center}
with \code{t\_total} the sum of the four times; $\chi$ the largest bond dimensions (\code{maxlinkdim}; for $\psi$ the largest of the two
subspace states); memories as stored elements $\times16$ bytes in MB ($10^6$ bytes), for $K$ and the sum over the two subspace states;
\code{rss\_MB = Sys.maxrss()/1e6}, the peak resident memory of the whole Julia process so far (it includes the Julia runtime and compiled
code, and, since $N$ increases along the scan, it is the peak up to and including this size); \code{start2} only the kind of start.
\end{function}

\paragraph{Driver and file logic.} If \code{RUN\_ZM = false} and the CSV exists, it is loaded with \code{readdlm(...; header=true)}. Otherwise
\code{zm\_run(4, 0)} is run once as a warm-up (compilation), then the scan runs with one line per finished size appended to
\code{data/zm\_progress.log} (flushed after each size); the CSV is written only at the end, header plus \code{writedlm} of all rows. This run is
not resumable: an interrupted scan leaves only the log. The cell prints one line per size (times, $\chi_H$, $\chi_K$, $\chi_\psi$, $K+\psi$ memory,
RSS, $E_f$, $\Delta Q$, gap, start of $\psi_2$).

Results: $\chi(K) = 31$--38 and $\chi(\psi) = 7$--26 over the whole scan, with no growth along the square scan; \code{t\_total} = 126 s and 140 s, RSS 1841 MB at both;
$K$ + MPS 1.40 MB at $2048\times2048$; the mirror start at every size; at $2048\times2048$: $E_f = 3.712285\times10^{-4}$, $\Delta Q = -3.446600\times10^{-3}$, gap 0.5529, the converged values from $L = 6$ on.

\subsection{Section 8: a substrate the zero modes leak into}

\code{RUN\_CAND = true}; \code{cand\_file = "tn\_candidate.csv"}; \code{const PC = (μ = -3.0, Δ = 0.6, J = 3.0)} ($t$ and $\alpha$ keep their
defaults); \code{cand\_scan} $= (7,6), (8,8), (9,9), (10,10), (11,11)$, i.e.\ $128\times64$ and $256^2$ to $2048^2$ sites; $\ell = 64$, \code{S\_z}.

\begin{function}{cand\_run(Lx, Ly; ℓ=64)}
As \code{zm\_run}, with \code{msh\_substrate(g, S\_z; PC...)}. Returns the 24 columns
\begin{center}\small
\code{Lx, Ly, ell, N, t\_H, t\_K, t\_dmrg, t\_pair, t\_total, chi\_K, chi\_psi, mem\_K\_MB, mem\_psi\_MB, rss\_MB, E\_f, dQ, E\_gap, E\_sub1, E\_sub2, x1, x2, start2, C\_perp, mirror\_perp}
\end{center}
i.e.\ in addition $\lVert H\psi_1\rVert$, $\lVert H\psi_2\rVert$ (the check that the second DMRG state is a zero mode) and the full
\code{start2} tuple with the norms of the orthogonalized $C\psi_1$ and mirror image; without \code{D} and \code{chi\_H}. \code{rss\_MB} is the
process peak so far, so it includes section 7.
\end{function}

\paragraph{Resume logic.} \code{done} is the set of $(L_x, L_y)$ already in the CSV. If \code{RUN\_CAND}, the header is written when
\code{done} is empty, and each missing size is run and appended as one line as soon as it finishes (\code{open(..., "a")}), with a progress
line \code{done Nx×Ny: t s}. An interrupted run therefore resumes at the next size. Finally the file is read and printed: substrate, $N$,
total time, $\chi_K$, $\chi_\psi$, RSS, $E_f$, $\Delta Q$, gap, $\lVert H\psi_2\rVert$, $x_1\dots x_2$, and the start with its two norms. The markdown
gives a run time of about 2.5 hours. Results (mirror start at every size; mirror\_perp $=1.000$; $\lVert H\psi_2\rVert = 9.06$--$9.08\times10^{-6}$, i.e.\ $E_f$):
\begin{center}\small\begin{tabular}{lrrrc}
substrate & total & $\chi_K$ & $\chi_\psi$ & \code{C\_perp}\\\hline
$128\times64$ & 516 s & 31 & 115 & 0.002\\
$256^2$ & 1340 s & 31 & 137 & 0.386\\
$512^2$ & 1536 s & 35 & 114 & 0.000\\
$1024^2$ & 2976 s & 34 & 135 & 0.317\\
$2048^2$ & 2556 s & 37 & 114 & 0.000\\
\end{tabular}\end{center}
$\chi_\psi$ is 114--115 where $\psi_1$ came out a single Majorana (\code{C\_perp} $\approx0$) and 135--137 where it is a mixture; this is why $1024^2$ takes longer than $2048^2$. At $2048\times2048$: $E_f = 9.047709\times10^{-6}$, $\Delta Q = 4.583431\times10^{-4}$, against the exact $9.047709\times10^{-6}$ and $4.583435\times10^{-4}$.

\clearpage
%======================================================================
\section{Part II --- Python notebook (\code{msh\_substrate\_python.ipynb})}
%======================================================================

The same model as a sparse CSR matrix (SciPy), with dense diagonalization where it fits and shift-invert Lanczos with a sparse direct
solver beyond. Packages: NumPy, \code{scipy.linalg}, \code{scipy.sparse}, \code{scipy.sparse.linalg} (\code{eigsh}, \code{splu},
\code{LinearOperator}), \code{mumps} (python-mumps, MUMPS 5.8 sequential, conda-forge), \code{multiprocessing}, \code{resource}, matplotlib.

\subsection{Setup}

The thread variables are set to 1 before NumPy is imported. \code{DATA\_DIR = "data"}. Pauli matrices \code{s0, sx, sy, sz} and the
Kronecker products (first factor Nambu $\tau$, second spin $\sigma$) \code{tzs0} $=\tau_z\sigma_0$, \code{txs0}, \code{t0sx}, \code{t0sy},
\code{t0sz}, \code{tzsx}, \code{tzsy}. \code{S\_z} and \code{S\_tilt} are the same textures as in Julia.

\begin{function}{in\_fresh\_process(fn, *args)}
\begin{Verbatim}
ctx = mp.get_context("fork"); q = ctx.Queue()
def target():
    try:
        out, status = fn(*args), "ok"
    except BaseException as e:
        out, status = None, f"{type(e).__name__}: {e}".replace(",", ";").strip()
    q.put((out, resource.getrusage(resource.RUSAGE_SELF).ru_maxrss / 1024, status))
\end{Verbatim}
Runs \code{fn(*args)} in a forked child process and returns \code{(result or None, peak RSS in MB, status)}. The peak is the child's
\code{ru\_maxrss} (kB on Linux, divided by 1024). Fork is used so that \code{fn} may be a closure or lambda (nothing is pickled except the
result). Every calculation that builds a large matrix runs this way, for two reasons: the kernel stays small, and every measurement starts
from the same baseline (a forked child inherits the kernel's resident memory, about 110 MB per the markdown, which would otherwise grow
from one calculation to the next). Errors in the child (including \code{MemoryError}) become the status string \code{"Type: message"}, with
commas replaced by semicolons so it can go into a CSV. The parent polls the queue with a 5 s timeout and checks \code{p.is\_alive()}, so a
child that dies without reporting (e.g.\ killed by the out-of-memory killer) gives \code{(None, nan, "process died (exit code ...)")}
instead of hanging.
\end{function}

\subsection{Section 1: Hamiltonian}

\begin{function}{shift(n)}
$n\times n$ CSR matrix with ones on the first subdiagonal, $\ket i\to\ket{i+1}$, open boundary.
\end{function}

\begin{function}{chain\_sites(Lx, Ly, ell)}
Returns the chain columns \code{np.arange(x0, x0 + ell)} with $x_0 = (N_x-\ell)/\!/2$, and the chain row $N_y/\!/2$; as \code{chain\_x},
\code{chain\_row} in Julia.
\end{function}

\begin{function}{ham(Lx, Ly, ell, S=S\_z, t=1.0, mu=-7.0, alpha=0.9, Delta=2.4, J=5.8)}
\begin{Verbatim}
Tx = sp.kron(sp.kron(Iy, shift(Nx)), hx)
Ty = sp.kron(sp.kron(shift(Ny), Ix), hy)
H = sp.kron(sp.identity(N), -mu * tzs0 + Delta * txs0) + Tx + Tx.conj().T + Ty + Ty.conj().T
\end{Verbatim}
Sparse $4N\times4N$ BdG matrix (CSR) in the basis $4n+c$, $n = i_yN_x+i_x$. With this index a site operator is \code{kron(A\_y, B\_x)} and the
full matrix \code{kron(site operator, 4×4 spinor matrix)}. On-site $-\mu\tau_z+\Delta\tau_x$; x bonds \code{kron(I\_y, shift\_x)}$\otimes h_x$,
which are open within each row (no wrap to the next row, unlike the Julia register shift, so no mask is needed); y bonds
\code{kron(shift\_y, I\_x)}$\otimes h_y$; $h_x = -t\tau_z+i\alpha\tau_z\sigma_y$, $h_y = -t\tau_z-i\alpha\tau_z\sigma_x$; both plus their conjugate
transposes. The exchange adds \code{kron(diag(v), t0s\_k)} with $v_n = JS_k(i_x)$ on the chain sites, skipped when $\max|v|\le10^{-14}$.
Used by every section.
\end{function}

\subsection{Section 2: check against the tensor-network Hamiltonian}

\begin{function}{julia\_to\_python(Hj, N)}
\begin{Verbatim}
W[4*n + 0, 0*N + n] = 1            # c_up      <- (nb = 0, up)
W[4*n + 1, 1*N + n] = 1            # c_dn      <- (nb = 0, dn)
W[4*n + 2, 2*N + 1*N + n] = 1      # c_dag_dn  <- (nb = 1, dn)
W[4*n + 3, 2*N + 0*N + n] = -1     # -c_dag_up <- (nb = 1, up)
\end{Verbatim}
Converts a dense Julia matrix (index $n_b\cdot2N+\sigma N+n$) to the Python basis: $W H_{\rm TN}W^T$ with the signed permutation $W$ (comments in
the listing transliterated). The minus sign is the $-c^\dagger_\uparrow$ of kwant's spinor.
\end{function}

\paragraph{Check.} For every key of \code{data/tn\_check\_H.npz} the cell parses $L_x$, $L_y$, $\ell$ and the texture from the key, builds
\code{ham(...)} with \code{S\_z} or \code{S\_tilt}, and prints $\max|WH_{\rm TN}W^T-H_{\rm py}|$ element by element; it then asserts that the
largest difference is below $10^{-10}$. Printed: largest difference $9.8\times10^{-14}$.

\subsection{Section 3: lowest levels against substrate width}

$L_x = 6$, $\ell = 48$, \code{S\_z}, $L_y = 0\dots6$: \code{eigsh(H, k=6, sigma=0, which="LM", return\_eigenvectors=False)} (shift-invert
with SciPy's default factorization). The sorted $|E|$ give the pair, \code{E[0]}, and the gap, \code{E[2]} (pairs of $\pm E$ appear twice). Run in
the kernel (largest $D = 16\,384$). The purpose stated in the markdown: the substrate at $\mu = -7$ is gapped, so the gap converges after a
few rows, and the converged values are a correctness check at sizes where dense diagonalization is no longer possible.
The gap is 0.2563 for one row, 0.4283 and 0.5613 for 2 and 4 rows, and converged at 0.5491 from 8 rows ($L_y = 3$) on; the pair $|E|$ is $2.6\times10^{-11}$ for one row and $\sim10^{-14}$ with a substrate.

\subsection{Section 4: sparse matrix size against system size}

\begin{function}{size\_info(Lx, Ly, ell)}
Builds \code{ham} and returns $(D, \text{nnz}, \text{CSR MB}, t_{\rm build})$, with the CSR memory \code{data + indices + indptr} bytes. Run in a
forked process for each size of the width scan $(6, L_y, 48)$ and the square scan $(L, L, 3\cdot2^L/4)$, $L = 4\dots10$ (the TN scan goes to
$L = 11$). At $1024\times1024$: $8.38\times10^7$ nonzeros, 1693 MB in CSR, built in 1.5 s.
\end{function}

\subsection{Section 5: zero-mode benchmark, dense and sparse}

Settings: \code{RUN\_ZM = True}; \code{DENSE\_MAX\_D = 16384} (dense only up to $64\times64$ sites); \code{zm\_file = "py\_zero\_modes.csv"};
\code{zm\_scan} as in Julia without $2048\times2048$ (13 sizes, increasing $N$); $\ell = 12$, \code{S\_z}.

\begin{function}{pair\_values(E, V)}
From a few eigenpairs around zero: the two smallest $|E|$ are the pair, and \code{ip} is the one with $E>0$. Returns
$E_f = E_{ip}$, $\Delta Q = \sum_n(|v_{n0}|^2+|v_{n1}|^2) - (|v_{n2}|^2+|v_{n3}|^2) = \bra{\psi_f}\tau_z\ket{\psi_f}$ (eigenvectors are normalized), and
the gap, the third smallest $|E|$.
\end{function}

\begin{function}{solve(Lx, Ly, method, ell=12)}
Zero modes of one system by one method; returns (eigensolver time, $E_f$, $\Delta Q$, gap). The time includes the factorization but not the
build of $H$.
\begin{itemize}
\item \code{"superlu"}: \code{splu(H.tocsc(), permc\_spec="MMD\_AT\_PLUS\_A")}, then
  \code{eigsh(H, k=4, sigma=0, which="LM", OPinv=LinearOperator(..., matvec=lu.solve))}, i.e.\ Lanczos on $H^{-1}$. The ordering is the
  minimum-degree ordering of $A^T+A$, suited to the structurally symmetric matrix; the markdown states that it halves the fill compared with
  \code{eigsh}'s default (COLAMD) and makes the solve 2.3--2.5 times faster ($256^2$: 21 s to 9 s; $512^2$: 276 s to 110 s). SuperLU reports a failed
  allocation by printing to stderr and raising an empty \code{MemoryError}; \code{solve} re-raises it as
  \code{MemoryError("SuperLU: not enough memory to perform factorization")} so the status in the CSV is informative. At $1024\times1024$
  this happens within a second with a few GB in use on a 188 GB machine (SuperLU's limits, not the machine's), and that size is recorded as
  failed.
\item \code{"mumps"}: \code{mumps.Context().factor(H.tocoo(), ordering="metis")} (analysis and numerical factorization with METIS nested
  dissection), then the same \code{eigsh} with \code{ctx.solve}. Per the markdown, METIS gave the smallest factorization at $256^2$ (2.6 s, against
  3.1 s for SCOTCH and 3.2 s for AMD), and MUMPS's 64-bit indices let it factorize $1024\times1024$.
\item \code{"dense"}: \code{la.eigh(H.toarray(), subset\_by\_index=[2N-2, 2N+1])}, the four eigenpairs in the middle of the spectrum, about twice
  as fast as the full spectrum; the $O(D^3)$ tridiagonal reduction dominates either way.
\end{itemize}
\end{function}

\paragraph{Driver and resumable CSV.} Columns \code{Lx, Ly, ell, N, D, method, t, rss\_MB, E\_f, dQ, E\_gap, status}. \code{done} is the set of
$(L_x, L_y, \text{method})$ already in the file. If \code{RUN\_ZM}: one warm-up \code{solve(4, 0, m)} per method in the kernel; the baseline
peak of a forked process is printed (\code{in\_fresh\_process(lambda: None)}); the header is written if nothing is done; then for each size and
method (dense only while $D\le$ \code{DENSE\_MAX\_D}) a missing row is computed with \code{in\_fresh\_process(solve, ...)} and appended at once. A
failed or killed run is recorded with NaN values and its status. Finally the file is read, sorted by $(N, \text{method})$ and printed.
\code{RUN\_ZM = False} only loads it; deleting the file starts over. MUMPS at $1024\times1024$: 1013 s and 28.6 GB (the slowdown of this point is explained in Part III, section 5); dense at $64\times64$: 886 s and 4.4 GB; SuperLU at $1024\times1024$: \code{MemoryError: SuperLU: not enough memory to perform factorization}.

\subsection{Section 6: MUMPS estimates beyond the benchmark}

\begin{function}{mumps\_estimate(L)}
Builds \code{ham(L, L, 12)} (COO), runs only MUMPS's analysis step (\code{ctx.analyze(H, ordering="metis")}: ordering and symbolic
factorization, no numerical work) and returns
$(L, L, N, D$, \code{est\_nonzeros}, \code{est\_flops}, \code{est\_mem\_incore}, \code{est\_mem\_ooc}, $t_{\rm analysis})$ from
\code{ctx.analysis\_stats}: the predicted number of entries in the factors, the floating-point operations of the factorization, and the
in-core and out-of-core memory in MB.
\end{function}
\code{RUN\_EST = True} (or a missing file) runs it in a forked process for $L = 7\dots11$ and rewrites \code{py\_mumps\_estimates.csv}; then the
table is printed. The markdown states that the in-core estimate is 7\,\% ($512^2$) and 2.5\,\% ($1024^2$) below the measured peak of section 5,
so the $2048\times2048$ value (about 120 GB; roughly two hours at the speed measured at $1024^2$) is reliable. The estimates depend only on
the sparsity pattern, which is the same at both parameter points with \code{S\_z}. At $2048\times2048$: $5.76\times10^9$ factor entries, $1.54\times10^{13}$ flops, 121184 MB in core, 14703 MB out of core.

\subsection{Section 7: charge difference, mask method and single-mode formula}

The CZ, QFT and period-finding notebooks and \code{msh\_chain\_python\_ed.ipynb} compute the charge difference from the many-body densities of
the two logical states (the \emph{mask method}); this notebook and the TN code use the single vector $\psi_f$. The section shows that they are
the same quantity.

\paragraph{Mask method.} $M$ is the full $4N\times4N$ unitary of eigenvectors sorted by energy: column $2N$ is $\psi_f$ ($+E_f$), column $2N-1$ its
partner $C\psi_f$ ($-E_f$), columns $0\dots2N-1$ the negative-energy states. Logical state $\ket0$ fills columns $0\dots2N-1$; $\ket1$ swaps column
$2N-1$ for column $2N$. With $G_j = M_{:,\text{mask}_j}M_{:,\text{mask}_j}^\dagger$,
\[
\rho_j(\mathbf r) = \sum_\sigma(G_j)_{\mathbf r\sigma,\mathbf r\sigma},\qquad \delta n = \rho_1-\rho_0,\qquad \Delta Q = \sum_{\mathbf r}\delta n(\mathbf r).
\]
\paragraph{Reduction to one vector.} The two masks share every column but one, so
\[
G_1 - G_0 = \ket{\psi_f}\bra{\psi_f} - \ket{C\psi_f}\bra{C\psi_f}.
\]
The particle part of $C\psi_f$ is the complex conjugate of the hole part $v$ of $\psi_f$ up to a spin rotation ($C = \tau_y\sigma_yK$), so the
particle diagonal is $\delta n(\mathbf r) = \sum_\sigma|u_\sigma(\mathbf r)|^2 - \sum_\sigma|v_\sigma(\mathbf r)|^2$ and
$\Delta Q = \bra{\psi_f}\tau_z\ket{\psi_f}$. The bulk drops out of the difference; the identity is exact, also during a ramp (all columns evolve
with the same unitary) and column by column for several modes. The mask method needs all of $M$ (dense \code{eigh}, $O(D^3)$ time, $O(D^2)$
memory; per the markdown $64\times64$ sites take 4.3 GB and 886 s, and $1024\times1024$ would need about $2.8\times10^5$ GB); its extra output, the
absolute densities $\rho_0,\rho_1$, is the bulk charge, which cancels in every compared quantity. Hence only the single-mode formula is used for the
scaling study. Conventions that must match: $\psi_f$ is the $+E_f$ member (the partner gives $-\Delta Q$), and the density counts electrons,
the components $(c_\uparrow, c_\downarrow)$.

\begin{function}{mask\_check(Lx, Ly, ell=12)}
\begin{Verbatim}
masks = np.zeros((2, 4*N), bool)
masks[1, 2*N] = True                                 # |1>: psi_f filled
masks[:, :2*N] = np.flip(~masks[:, 2*N:], axis=1)    # column 2N-1-k = C(column 2N+k)
\end{Verbatim}
One full dense \code{la.eigh} (ascending energies, timed). The masks are built as in \code{msh\_chain\_python\_ed.ipynb}: the upper half
says which positive-energy columns are filled, and the lower half is its flipped complement, because column $2N-1-k$ is the partner of column
$2N+k$. For each mask $\rho_j$ is the particle diagonal of $G_j$, computed as the row sums of $|M_{:,\text{mask}}|^2$ over the particle
components. The single-mode density uses column $2N$ alone. $C\psi_f$ is computed as \code{kron(sy, sy)} applied to the conjugated per-site
spinors ($C = \tau_y\sigma_yK$) and compared with column $2N-1$. Returns a dict with $N$, $D$, $E_f$, the \code{eigh} time, the overlap
$|\langle M_{:,2N-1}|C\psi_f\rangle|$, $\Delta Q$ from both methods, $\max|\delta n_{\rm mask}-\delta n_{\rm single}|$ and both densities as
$N_y\times N_x$ arrays.
\end{function}
Run in forked processes for $64\times1$, $64\times4$, $16\times16$, $32\times32$ (\code{mask\_scan}); the cell prints one line per system and plots
$\delta n(i_x,i_y)$ of the $32\times32$ system (mask method) and $|\delta n_{\rm mask}-\delta n_{\rm single}|$. Overlap $1.000000000000$ in all four cases, $|\Delta Q_{\rm mask}-\Delta Q_{\rm single}|\le2.1\times10^{-15}$, $\delta n$ differences $\le2.0\times10^{-12}$.

\subsection{Section 8: a substrate the zero modes leak into}

\subsubsection*{8a.\ \texorpdfstring{$\mu\times J$}{mu x J} scan for three values of \texorpdfstring{$\Delta$}{Delta}}

Fixed geometry $64\times32$ ($L_x = 6$, $L_y = 5$), $\ell = 32$ in row 16. Module-level masks over the $(i_y, i_x)$ grid: \code{ends\_s}, within 3
columns of either chain end (all rows); \code{mid\_s}, the middle third of the chain columns; \code{dy\_s} $=|i_y-i_{y0}|$.

\begin{function}{scan\_point(args)}
For \code{args = (Delta, mu, J)}: \code{eigsh(ham(...), k=6, sigma=0, which="LM")}, levels sorted by $|E|$. With $\rho$ the total weight
(all four components) per site of the lowest state, returns
$(\Delta, \mu, J, E_0, E_2, E_4$, \code{w\_ends}, \code{w\_mid}, \code{w\_row}, \code{rms\_dy}, \code{w\_far}$)$: the first three distinct levels
$|E|$ (indices 0, 2, 4 of the $\pm$ pairs), the weight near the chain ends, in the middle third, on the chain row, the rms distance
$\sqrt{\sum\rho\,\Delta y^2}$ from the chain row, and the weight four or more rows away.
\end{function}
Grid: $\Delta\in\{2.4, 1.2, 0.6\}$, $\mu = -7\dots0$ in steps of 0.25, $J = 1\dots8$ in steps of 0.5: $3\times29\times15 = 1305$ points,
\code{Pool(8).map(..., chunksize=4)}, written to \code{py\_mu\_scan.csv} (\code{RUN\_MU = False} loads it). \textbf{Topological criterion:}
$E_0/E_2<0.05$ (a pair well below the next level) and \code{w\_ends} $>0.5$ (the lowest state sits at the chain ends). Figure: top row
$\log_{10}(E_0/E_2)$ per $\Delta$ with the topological region outlined in red; bottom row \code{rms\_dy} inside it; dashed line at the band
bottom $\mu = -4.741$. Printed: the number of topological points per $\Delta$, the topological point with the largest $E_2$ below and inside
the band, and four selected points, including $(0.6, -3, 3)$. 1305 points in 198 s; topological: 97 ($\Delta = 2.4$), 113 (1.2), 90 (0.6).

\subsubsection*{8b.\ The point \texorpdfstring{$\Delta = 0.6$, $\mu = -3$, $J = 3$}{Delta = 0.6, mu = -3, J = 3}}

\begin{function}{pair\_levels(Lx, Ly, ell)}
\code{eigsh} (6 levels) at \code{PC = dict(mu=-3.0, Delta=0.6, J=3.0)}. Returns $|E_0|$, $|E_2|$, $\Delta Q$ of the $+E$ member of the pair, and
the weight per site of that state as an $N_y\times N_x$ array.
\end{function}
Three checks, each size in a forked process: (1) splitting against chain length, $\ell = 16\dots128$ on $256\times32$, printing $E_0$, $E_2$,
$E_0/E_2$, $\Delta Q$; (2) convergence with substrate height, $N_x = 256$, $\ell = 64$, $N_y = 8\dots128$; (3) the weight per row, averaged over
$\pm\Delta y$, at $\Delta y = 0,1,2,4,\dots,48$ on $256\times128$. These produce the numbers quoted in the model description (decay length about 7
sites, $\ell = 64$, $e^{-|\Delta y|/2}$ decay). $E_0/E_2 = 1.0\times10^{-4}$ at $\ell = 64$ ($256\times32$); $E_0$, gap and $\Delta Q$ converge from $N_y = 64$ on ($E_0 = 9.0477\times10^{-6}$, gap 0.09025, $\Delta Q = 4.5834\times10^{-4}$); the row weight falls from 0.32 on the chain row to $3.6\times10^{-3}$ at $\Delta y = 8$ and $1.4\times10^{-11}$ at 48.

\subsubsection*{8c.\ MUMPS reference at the TN sizes}

\begin{function}{solve\_candidate(Lx, Ly, ell=64)}
\code{solve(..., "mumps")} at \code{PC}: METIS factorization and shift-invert \code{eigsh(k=4)}; returns (time including the factorization,
$E_f$, $\Delta Q$, gap) via \code{pair\_values}.
\end{function}
\code{cand\_scan} $= (7,6), (8,8), (9,9), (10,10)$, i.e.\ the TN sizes up to $1024\times1024$. Same resumable logic as section 5 with
\code{py\_candidate.csv} (columns \code{Lx, Ly, ell, N, t, rss\_MB, E\_f, dQ, E\_gap, status}); there is no \code{RUN} flag, missing sizes are
always run. The factorization depends only on the sparsity pattern, so the estimates of section 6 hold; the cell prints the table and a
$2048\times2048$ line with time $t_{1024}\times\text{flops}_{2048}/\text{flops}_{1024}$ and the in-core memory estimate, both marked
``\textasciitilde''. MUMPS: 1.0 s ($128\times64$), 10.7 s ($256^2$), 53.7 s ($512^2$), 295.1 s ($1024^2$); estimate at $2048^2$: $\sim$2326 s and $\sim$121 GB.

\clearpage
%======================================================================
\section{Part III --- Scaling notebook (\code{msh\_substrate\_scaling.ipynb})}
%======================================================================

Loads the CSV files of the other two notebooks and compares the cost of the methods; it computes nothing expensive. Inputs:
\code{tn\_zero\_modes.csv}, \code{py\_zero\_modes.csv}, \code{py\_mumps\_estimates.csv}, and in section 6 \code{tn\_candidate.csv},
\code{py\_candidate.csv}.

\subsection{Loading and helpers (cell 1)}

\code{tn} is the TN table sorted by $N$; \code{py} the Python table (with string columns); \code{est} the MUMPS estimates. Rows of \code{py} whose
status is not \code{"ok"} are printed as failed and left out of \code{M}, the dictionary of successful rows per method, sorted by $N$.

\begin{function}{key(r), ref(a, col)}
\code{key} returns $(L_x, L_y)$ of a row. \code{REF} maps sizes to the exact reference row: SuperLU rows, overwritten by MUMPS rows where MUMPS ran
(i.e.\ MUMPS wherever available). \code{ref(a, col)} returns the reference value of column \code{col} for each row of \code{a}, NaN where there is
none ($2048\times2048$).
\end{function}

\begin{function}{scan(a, which)}
Rows of the width scan (\code{Lx == 6}) or the square scan (\code{Lx == Ly}), sorted by $N$. $64\times64$ belongs to both.
\end{function}

\begin{function}{fit(N, t, npts)}
Power law $t = cN^p$ by least squares in $\log t$ against $\log N$ through the \code{npts} largest finite points; returns $(p, c, N_{\rm fit})$.
\end{function}

\paragraph{MUMPS at $2048\times2048$.} \code{t\_mumps\_2048} $= t_{\rm MUMPS}(1024^2)\cdot\text{flops}_{2048}/\text{flops}_{1024}$, the measured
$1024\times1024$ time scaled by the ratio of the estimated factorization flops. The cell prints the number of sizes per method and the estimate
in GB and hours. Printed: 121 GB in core, $\sim$7983 s (2.2 h).

\subsection{Section 1: time}

Two log--log panels (width scan, square scan) of the wall time against $N$: TN \code{t\_total} (build of $H$ and $K$, three DMRGs, pairing),
and the eigensolver time (including the factorization) of SuperLU, MUMPS and dense. On the square panel, dashed power-law fits through the
largest points (TN and MUMPS four, SuperLU and dense three) with the exponent in the legend, and the MUMPS estimate as an open diamond.
Square scan: TN (DMRG) $N^{0.14}$, SuperLU $N^{1.54}$, MUMPS $N^{1.51}$, dense $N^{3.07}$.

\subsection{Section 2: memory}

One log--log panel of peak memory against $N$: the Julia process RSS (whole process up to that size, dominated by runtime and compiled code), the
TN data itself (\code{mem\_K\_MB + mem\_psi\_MB}), the RSS of each Python method (its own forked process, baseline about 110 MB included), the
MUMPS in-core estimates (open diamonds), the dense-matrix line $16D^2$ bytes, and the 188 GB of the machine. Printed: the growth exponents of the
SuperLU and MUMPS RSS over their three largest square sizes (\code{fit}), measured against estimated MUMPS memory for $N\ge65\,536$, and the TN
data and process RSS at $2048\times2048$. RSS grows as $N^{0.98}$ (SuperLU) and $N^{1.05}$ (MUMPS) over the three largest sizes; MUMPS measured/estimated: 1552/1326 MB ($256^2$), 6594/6117 MB ($512^2$), 28636/27913 MB ($1024^2$), i.e.\ the measured RSS is 17\,\%, 8\,\% and 2.6\,\% above the estimate; TN data 1.40 MB and Julia RSS 1841 MB at $2048\times2048$.

\subsection{Section 3: accuracy}

$|E_f-E_f^{\rm exact}|$ and $|\Delta Q-\Delta Q^{\rm exact}|$ against $N$ for TN, SuperLU and dense, with \code{ref} as the exact value; differences
are clipped at $10^{-17}$ for the log scale (so SuperLU at sizes where it is its own reference shows at the floor). The TN values use
\code{cutoff\_K = 1e-20}. Printed: the largest TN errors over the sizes with a reference, the MUMPS values at $1024\times1024$, and the TN values at
$2048\times2048$ for comparison with the converged ones. Largest TN errors $2.5\times10^{-8}$ ($\Delta Q$) and $1.1\times10^{-11}$ ($E_f$); MUMPS at $1024^2$: $E_f = 3.7122847509\times10^{-4}$, $\Delta Q = -3.4466003632\times10^{-3}$; TN at $2048^2$: $3.7122846056\times10^{-4}$, $-3.4465999587\times10^{-3}$.

\subsection{Section 4: bond dimensions}

$\chi(H)$, $\chi(K)$ and the largest $\chi(\psi)$ against $N$ (semilog). The point of the plot, per the markdown: if they stay bounded, the TN cost
grows only through the number of qubits, $\log_2N+2$. Both stay bounded: $\chi(H) = 7$, $\chi(K) = 31$--38, $\chi(\psi)\le26$ at every size.

\subsection{Section 5: summary and crossover}

A table per TN size: times of TN, SuperLU, MUMPS, dense, the ratio MUMPS/TN (above 1: TN faster), and the four peak memories. The helper
\code{cell(m, k, col, fmt, w)} formats one entry: a dash if the method has no row for that size, \code{failed} if its status is not
\code{"ok"}. At $2048\times2048$ the MUMPS columns hold the estimates, marked \textasciitilde.

\paragraph{Crossover.} On the square scan, the points $(N, \log(t_{\rm MUMPS}/t_{\rm TN}))$ at the sizes where both ran (measured MUMPS only); between
the first pair of consecutive sizes where the log-ratio changes from negative to non-negative, $N_c$ is found by linear interpolation of the
log-ratio in $\log N$, and printed with $\sqrt{N_c}$ and the bracketing sizes. Printed: $N_c\approx3.47\times10^5$ sites ($\approx589\times589$), between $512^2$ and $1024^2$. The markdown adds that the MUMPS times at this point are inflated by subnormal floating-point arithmetic: with flush-to-zero the $1024^2$ factorization takes 107 s instead of 771 s, and the crossover would move to roughly $700^2$.

\paragraph{Caveat on the MUMPS times at the benchmark point} (markdown). At $1024\times1024$ MUMPS needs 1013 s (a separate rerun: 1074 s, of
which 902 s factorization), against 308 s at the second point with the same fill and the same peak memory (28.6 GB), and 14.5 times its
$512\times512$ time where $N^{1.5}$ would give about 8. Since the factorization does the same work on the same pattern, the difference is in the
arithmetic. A possible cause, not yet tested, is subnormal floating-point numbers: at $\mu = -7$ the Green's function decays within a fraction of
a site, factor entries can underflow to $\sim10^{-300}$, and arithmetic on subnormals is slow on x86. If so, the crossover is partly set by this
effect rather than by the cost of the algorithm. The $2048\times2048$ estimate inherits it, since it scales the measured $1024^2$ time.

\subsection{Section 6: a substrate the zero modes leak into}

Loads \code{tn\_candidate.csv} and \code{py\_candidate.csv} (successful rows only), sorted by $N$. The MUMPS time at $2048\times2048$ is estimated
as above from the measured $1024^2$ time at this point (\code{t\_mc\_2048}); the memory estimate is the same as at the benchmark point (same
pattern). The exact values \code{ex = (E\_f = 9.04770870e-06, dQ = 4.5834351e-04)} (converged MUMPS values of Python section 8c) are hard-coded.
A table per TN size: TN and MUMPS times and their ratio, TN process RSS and MUMPS RSS (estimates marked \textasciitilde), $\chi_\psi$, and the TN
errors $|E_f-\text{exact}|$, $|\Delta Q-\text{exact}|$. Two log--log panels (time, peak memory) show both methods at this point, the MUMPS estimate
as a dotted segment from $1024^2$ to $2048^2$, and the square-scan curves of the benchmark point as open symbols. The crossover is computed as in
section 5, here including the $2048^2$ estimate; if MUMPS is faster at every size the cell says so. Note that the TN RSS is the process peak
including section 7, while the MUMPS RSS is per forked process.
MUMPS/TN $= 0.00, 0.01, 0.03, 0.10, 0.91$ from $128\times64$ to $2048^2$; TN errors $\le2.6\times10^{-12}$ ($E_f$) and $\le9.0\times10^{-10}$ ($\Delta Q$); printed: no crossover up to $2048\times2048$ (MUMPS is faster at every size, with its estimate at $2048^2$), but the TN needs 2.3 GB there against $\sim$121 GB.

\end{document}
